% GL1F formal-results fragment.
% Requires: amsmath, amssymb, amsthm, booktabs.
% Expected theorem environments: definition, lemma, theorem, proposition,
% corollary, remark.

\section{Formal model and correctness results}
\label{sec:formal-results}

This section separates three objects that are easy to conflate:
(i) the mathematical forest, (ii) its GL1F byte representation, and
(iii) a registry entry that points to EVM code chunks.  The equivalence result
for the first two is exact under the representation, input, depth, and
arithmetic premises below.  Its application to a registry entry is also
conditional on a well-formed storage manifest, because the deployed registry
does not itself recompute the model hash or validate every serialization
invariant.

\subsection{Canonical representation}

Let $\mathbb{I}_{32}=[-2^{31},2^{31}-1]\cap\mathbb{Z}$.  A complete binary
tree of depth $d$ contains
\begin{equation}
  L(d)=2^d
  \quad\text{leaves and}\quad
  I(d)=2^d-1
  \quad\text{internal nodes}.
\end{equation}
An internal node occupies eight bytes: a little-endian unsigned 16-bit feature
index, a little-endian signed 32-bit threshold, and two reserved bytes.
A leaf occupies four bytes containing a little-endian signed 32-bit value.
Consequently, one serialized tree occupies
\begin{equation}
  b(d)=8I(d)+4L(d)=12\cdot 2^d-8
  \label{eq:tree-bytes}
\end{equation}
bytes.

\begin{definition}[Canonical GL1F v1 model]
\label{def:gl1f-v1}
A canonical publication-profile scalar model is a byte string $B$ consisting
of a 24-byte header and $T$ trees, where $F\ge1$, $1\le d\le12$, and
$0\le T\le2^{32}-1$.  A zero-tree v1 object is the well-defined base-only
case.  Its header contains magic \texttt{GL1F}, version one, feature count
$F$, depth $d$, tree count $T$, base value
$a_0\in\mathbb{I}_{32}$, and scale
$s\in[1,2^{32}-1]\cap\mathbb{Z}$.  Reserved bytes are zero, every node feature
is in $\{0,\ldots,F-1\}$, and
\begin{equation}
  |B|=24+T b(d).
  \label{eq:v1-length}
\end{equation}
For a deployed-registry EVM application we additionally require equality
between serialized and registered $F,d,T,a_0,$ and $s$, and
$T\le2^{16}-1$.  The latter is a deployed-interface restriction: the v1 wire
field for $T$ is unsigned 32-bit, whereas the deployed registry stores total
tree count as unsigned 16-bit.
\end{definition}

\begin{definition}[Canonical GL1F v2 model]
\label{def:gl1f-v2}
A canonical publication-profile vector model has $F\ge1$, $1\le d\le12$,
$K\ge2$ outputs, $R\ge1$ trees per output,
base vector $(a_{0,0},\ldots,a_{0,K-1})\in\mathbb{I}_{32}^K$, and output-major
tree order.  Its exact length is
\begin{equation}
  |B|=24+4K+KRb(d).
  \label{eq:v2-length}
\end{equation}
The header has magic \texttt{GL1F}, version two, positive $F$ and $s$, and
zero reserved fields.  Every node feature is smaller than $F$.  Registered
feature count and scale equal the serialized values.  Publication through the
deployed registry additionally requires $KR\le2^{16}-1$, because
that interface stores the total tree count in a 16-bit field; this is narrower
than the core format's separate 16-bit $K$ and 32-bit $R$ fields.
\end{definition}

The format does not contain a task discriminator.  Identical v1 bytes can be
described externally as regression or binary-logit inference, and identical
v2 bytes can be described as multiclass or multilabel inference.  Here
$\operatorname{keccak256}$ denotes Ethereum's legacy Keccak-256 function, not
the standardized FIPS SHA3-256 variant.  Under the usual collision- and
second-preimage-resistance assumptions, $\operatorname{keccak256}(B)$ binds a
practical identifier to the exact serialized bytes $B$.  It is not a semantic
normal form: distinct canonical cores can denote the same extensional integer
function.  The digest also does not bind a task label, feature semantics,
preprocessing contract, dataset, or intended use.  Those claims require a
separately authenticated semantic manifest.

The explicit depth premise matches the current first-party interface.  It is
substantially narrower than the 16-bit wire field and prevents an apparently
valid header from requesting an impractically large complete tree or an
evaluator-specific shift result.

\subsection{Chunked EVM representation}

Let $c$ be the nominal payload of every nonfinal chunk and therefore the
logical byte-offset stride; the final payload may be shorter.  Let
$N=\lceil |B|/c\rceil$.  The canonical publisher deploys each payload through
\texttt{ModelStore.write}, creating runtime code
\[
  C_j=\texttt{GL1C}\,\Vert\,
  B[jc:\min\{(j+1)c,|B|\}].
\]
It then creates a table whose data are $N$ 32-byte words; the low 20 bytes of
word $j$ are the address of $C_j$.  The table is itself prefixed by
\texttt{GL1C}.

\begin{definition}[Well-formed storage manifest]
\label{def:manifest}
For an EVM state $\sigma$ and registry record $r$ in that state, a manifest
$(P,c,N,\ell)$ is well formed for $B$ at $(\sigma,r)$ when:
\begin{enumerate}
  \item $P\ne0$, $4\le c\le24572$, $1\le N\le767$,
        $\ell=|B|$, and $N=\lceil\ell/c\rceil$;
  \item table address $P$ has runtime code in $\sigma$ of exact length $4+32N$, beginning
        with \texttt{GL1C};
  \item its $j$th word has zero high 12 bytes and resolves through its low
        20 bytes to nonzero address $C_j$;
  \item every nonfinal $C_j$ in $\sigma$ has exact code
        $\texttt{GL1C}\Vert B[jc:(j+1)c]$, and the final chunk has the exact
        code $\texttt{GL1C}\Vert B[(N-1)c:\ell]$; and
  \item the corresponding fields of $r$ equal $P,c,N,\ell$.
\end{enumerate}
\end{definition}

The lower bound $c\ge4$ is semantic, not cosmetic.  Runtime primitive fields
are at most four bytes, while the implemented boundary reader consults at most
two chunks.  A field could span more than two logical chunks if $c<4$.

\begin{lemma}[Pointer capacity]
\label{lem:pointer-capacity}
A pointer table created by the deployed storage contract contains at most
\[
  \left\lfloor \frac{24572}{32}\right\rfloor=767
\]
chunk addresses.
\end{lemma}

\begin{proof}
\texttt{ModelStore.write} accepts at most 24572 data bytes, and each table slot
occupies exactly 32 data bytes.  The four-byte \texttt{GL1C} prefix is added by
the store and is not part of the payload limit.
\end{proof}

For a manifest satisfying Definition~\ref{def:manifest}, define the specified
boundary reader $\operatorname{Read}_{\sigma,M}(o,n)$ to compute
$j=\lfloor o/c\rfloor$ and $r_o=o\bmod c$, copy from code offset $4+r_o$ in
chunk $j$, and, only when $r_o+n>c$, concatenate the remaining bytes from code
offset four in chunk $j+1$.

\begin{lemma}[Exact byte reconstruction]
\label{lem:reconstruction}
Given Definition~\ref{def:manifest}, every runtime primitive read of
$n\in\{1,2,4\}$ bytes at a model offset $o\in\mathbb Z_{\ge0}$ satisfying
$o+n\le |B|$ returns exactly $B[o:o+n]$.
\end{lemma}

\begin{proof}
By Euclidean division, write $o=jc+r$ with
$j=\lfloor o/c\rfloor$ and $0\le r<c$.  If $r+n\le c$, the runtime resolves
slot $j$ and copies from code offset $4+r$, which is exactly
$B[o:o+n]$.  Otherwise it copies $c-r$ bytes from chunk $j$ and
$n-(c-r)$ bytes from the start of chunk $j+1$.  Since $n\le4\le c$, no third
chunk is required.  The exact chunk-length and ordering clauses in
Definition~\ref{def:manifest} make the concatenation equal to $B[o:o+n]$.
\end{proof}

\begin{remark}
The \texttt{GL1C} check detects zero and many accidental pointers, but it is
not content authentication.  Arbitrary EVM code beginning with the same four
bytes passes that test.  Exact reconstruction is therefore a property of the
manifest conditions, not of the magic prefix alone.
\end{remark}

\subsection{Traversal and termination}

Let $i_j$ be the heap index after $j$ decisions, with $i_0=0$.  For packed
feature $x_q$ and threshold $\theta_q$, a node updates
\begin{equation}
 i_{j+1} =
 \begin{cases}
  2i_j+2, & x_q>\theta_q,\\
  2i_j+1, & x_q\le\theta_q.
 \end{cases}
 \label{eq:index-update}
\end{equation}

\begin{lemma}[Traversal interval]
\label{lem:traversal-interval}
After $j$ decisions,
\begin{equation}
  2^j-1\le i_j\le2^{j+1}-2.
  \label{eq:index-interval}
\end{equation}
\end{lemma}

\begin{proof}
For $j=0$, both bounds equal zero.  Assume the claim at $j$.
The smallest possible child is
$2(2^j-1)+1=2^{j+1}-1$, and the largest is
$2(2^{j+1}-2)+2=2^{j+2}-2$.  These are precisely the bounds for $j+1$.
\end{proof}

\begin{corollary}[Leaf-index safety]
\label{cor:leaf-safety}
After exactly $d$ decisions,
\[
  0\le i_d-I(d)\le L(d)-1.
\]
Thus a traversal of a well-formed tree selects exactly one serialized leaf and
cannot address outside that tree's leaf array.
\end{corollary}

\begin{proof}
Substitute $j=d$ in Lemma~\ref{lem:traversal-interval} and subtract
$I(d)=2^d-1$.
\end{proof}

\begin{lemma}[Forced-subtree inertness]
\label{lem:forced-subtree}
Suppose a terminal node denoting signed int32 value $v$ is replaced by a
complete forced continuation whose remaining internal nodes use feature zero
and threshold $2^{31}-1$, and whose descendant leaves all contain that same
$v$. For every $q\in\mathbb{I}_{32}^F$, traversal through that continuation
returns $v$; hence the replacement preserves the terminal value over the
entire packed-input domain.
\end{lemma}

\begin{proof}
Definitions~\ref{def:gl1f-v1} and \ref{def:gl1f-v2} give $F\ge1$, so feature
zero exists.  Because
$q_0\le2^{31}-1$, the strict predicate
$q_0>2^{31}-1$ is false at every inserted node and traversal follows the left
child at each remaining level.  The reached descendant leaf contains $v$.
Moreover, assigning $v$ to every descendant leaf makes the value independent
even of the particular descendant selected.  Thus the forced continuation
cannot change the tree output for any int32 input vector.
\end{proof}

\begin{theorem}[Finite work and logical traversal state]
\label{thm:complexity}
For a canonical v1 model, inference performs exactly $Td$ node decisions and
$T$ leaf reads.  For a canonical v2 model, inference performs exactly $KRd$
node decisions and $KR$ leaf reads.  An abstract streaming evaluator that
does not buffer the model needs $O(1)$ scalar logical traversal state and
$O(K)$ vector logical traversal state, excluding input storage and
implementation-specific read buffers.  Model space is $24+Tb(d)$ for v1 and
$24+4K+KRb(d)$ for v2.
\end{theorem}

\begin{proof}
Each complete tree executes the loop in Equation~\ref{eq:index-update} exactly
$d$ times, followed by one leaf read.  The scalar outer loop visits $T$ trees.
The vector loops visit $R$ trees for each of $K$ outputs.  No inference loop is
recursive or conditioned on a model value other than the bounded header
counts.  A streaming evaluator stores one scalar accumulator or a length-$K$
result vector together with a constant number of indices and decoded
primitive values.  The model-space formulae follow from
Equation~\ref{eq:tree-bytes}.
\end{proof}

\begin{remark}[Bounded does not imply affordable]
The theorem establishes a finite operation count.  It does not guarantee that
an arbitrary encoded dimension fits a transaction gas limit.  The publication
profile and empirical gas envelope must be reported separately.  The deployed
Solidity boundary-reader allocates a fresh dynamic byte buffer when one
primitive spans two chunks.  Since EVM memory allocation is monotonic within a
call, its transient memory may grow linearly with the number of such boundary
reads even though every primitive read is at most 32 bytes.  This is an
implementation cost, not additional logical traversal state.
\end{remark}

\subsection{Cross-runtime equivalence}

A \emph{conforming forest evaluator} uses the specified little-endian
signed/unsigned decoding and strict comparison, accumulates the mathematical
integer sums below without overflow or loss of precision, and selects the
lowest index among equal maximum logits when returning a class. A conforming
local evaluator first validates the version-specific canonical definition. A
\emph{conforming EVM evaluator} at $(\sigma,r)$ obtains every core primitive
through $\operatorname{Read}_{\sigma,M}$, uses the version-aware registry
relations below, and otherwise applies the same forest semantics. These are
specification definitions; source inspection and finite compiled/deployed
tests separately assess the deployed \texttt{ForestRuntime} path. Client-side
preprocessing may construct the packed integer vector, but real-number
quantization is not part of the theorem. Applying the theorem to an external
endpoint additionally requires its enabled, authorization, and fee conditions
to reach the internal evaluator with the packed input unchanged.

For a scalar model, define the integer function
\begin{equation}
  f_B(q)=a_0+\sum_{t=0}^{T-1}v_{t,\pi_t(q)},
  \label{eq:scalar-function}
\end{equation}
where $\pi_t(q)$ is the leaf selected by
Equation~\ref{eq:index-update}.  For a vector model,
\begin{equation}
  f_{B,k}(q)=a_{0,k}+\sum_{t=0}^{R-1}v_{k,t,\pi_{k,t}(q)}.
  \label{eq:vector-function}
\end{equation}

\begin{theorem}[Specification-level conforming-evaluator equivalence]
\label{thm:equivalence}
Let $B$ be canonical under Definition~\ref{def:gl1f-v1} or
Definition~\ref{def:gl1f-v2}, let the EVM manifest satisfy
Definition~\ref{def:manifest} at the same $(\sigma,r)$ used for evaluation, and
require the following nine version-aware relations for $r$
\[
\begin{gathered}
nFeatures=F,\qquad
nTrees=\begin{cases}T&\text{(v1)},\\KR&\text{(v2)},\end{cases}
\qquad depth=d,\\
baseQ=\begin{cases}a_0&\text{(v1)},\\0&\text{(v2 registry convention)},\end{cases}
\qquad scaleQ=s,\qquad tablePtr=P,\\
totalBytes=\ell,\qquad numChunks=N,\qquad chunkSize=c.
\end{gathered}
\]
Let every evaluator receive the identical $4F$-byte sequence formed by packing
$q\in\mathbb{I}_{32}^F$ as signed little-endian int32 values. Then conforming
local and EVM evaluators return $f_B(q)$ for v1 and
$(f_{B,k}(q))_{k=0}^{K-1}$ for v2. Their v2 class operation returns
\[
 \left(\min\operatorname*{arg\,max}_{0\le k<K} f_{B,k}(q),
       \max_{0\le k<K} f_{B,k}(q)\right),
\]
where the minimum implements the specified lowest-index tie rule.
\end{theorem}

\begin{proof}
Lemma~\ref{lem:reconstruction} gives identical serialized bytes.  The v2 EVM
and local evaluators decode the same header, node, threshold, and leaf integers.
The conforming v1 EVM evaluator uses registered inference fields rather than parsing
the header; the theorem's nine-relation premise makes those fields
identical to the local decoder's values, while the node, threshold, and leaf
integers still follow from Lemma~\ref{lem:reconstruction}.  At the root, every
evaluator compares the same $q_f$ and $\theta_q$ using the same strict
predicate.  If their indices agree before a level, the identical predicate and
Equation~\ref{eq:index-update} give the same next index.  Induction and
Corollary~\ref{cor:leaf-safety} give the same selected leaf for every tree.
Finally, Equations~\ref{eq:scalar-function} and
\ref{eq:vector-function} add the same finite integer sequence to the same base
term, with the exact-accumulation premise preserving that mathematical sum.
Applying the same lowest-index argmax scan to the identical v2 vector gives
the same class index and maximum integer logit.
\end{proof}

\begin{remark}
This handwritten theorem concerns conforming evaluator specifications in one
state and record. Source inspection supports the mapping from the deployed
Solidity path to that specification; compiled tests and pinned observations
supply finite implementation evidence. None is a mechanized EVM-bytecode proof,
verified Solidity-to-bytecode refinement, or source-to-live-bytecode provenance proof.
\end{remark}

\begin{proposition}[Exact JavaScript accumulation under the deployment profile]
\label{prop:js-safe}
All integer accumulators permitted by the canonical ModelStore/registry
profile are strictly within JavaScript's exact integer range.
\end{proposition}

\begin{proof}
For v1, the registry's 16-bit tree count gives
\[
 |f_B(q)|\le (65535+1)2^{31}=2^{47}<2^{53}.
\]
For v2 with $d\ge1$, Equation~\ref{eq:tree-bytes} gives $b(d)\ge16$.
Lemma~\ref{lem:pointer-capacity} bounds total core payload by
$767\cdot24572=18846724$ bytes.  Since $K\ge2$,
\[
 R < \frac{18846724}{2\cdot16}<589000.
\]
Hence every output and every partial sum obeys
\[
 |f_{B,k}(q)|<(589001)2^{31}<2^{51}<2^{53}.
\]
All int32 inputs and summands are themselves exactly representable.
\end{proof}

\begin{remark}
The theorem concerns inference from fixed bytes.  Byte-identical training is a
separate empirical claim because split scoring, gradients, and early stopping
use floating-point arithmetic.  It requires fixed preprocessing, operation
order, feature/class order, seeded random-number consumption, and compiler
flags, plus a parity test matrix.
\end{remark}

\subsection{Fixed-point approximation}

For finite $z\in\mathbb{R}$ and integer $s>0$, define the abstract exact-real
quantizer with JavaScript's half-toward-positive-infinity tie rule
\begin{equation}
  Q_s(z)=\operatorname{clamp}_{\mathbb{I}_{32}}
  \left(\left\lfloor sz+\frac12\right\rfloor\right).
  \label{eq:quantizer}
\end{equation}
This definition includes the asymmetric negative-half convention of
\texttt{Math.round}; for example, $Q_1(-1.5)=-1$.

This is a mathematical abstraction.  Production JavaScript first computes the
binary64 product
$\operatorname{fl}(sz)=sz+\delta_s(z)$ and then applies
\texttt{Math.round}.  The two integers agree when that multiplication
perturbation does not cross an integer rounding boundary.  Without that
premise, and provided the rounded production result remains in the int32
range, the direct implementation bound is
\[
 \left|
 \frac{\operatorname{Math.round}(\operatorname{fl}(sz))}{s}-z
 \right|
 \le \frac{1/2+|\delta_s(z)|}{s},
\]
which can be slightly larger than $1/(2s)$.  If int32 saturation occurs,
neither bound applies.

\begin{lemma}[Abstract single-value quantization error]
\label{lem:quant-error}
If Equation~\ref{eq:quantizer} does not saturate, then
\[
 \left|\frac{Q_s(z)}{s}-z\right|\le\frac{1}{2s}.
\]
\end{lemma}

\begin{proof}
For $u=sz$, the integer $\lfloor u+1/2\rfloor$ lies in
$[u-1/2,u+1/2]$.  Divide by positive $s$.
\end{proof}

\begin{theorem}[Stable-path output error]
\label{thm:stable-path-error}
Consider an ideal scalar forest whose base and the $T$ leaves selected for
input $x$ are independently quantized by Equation~\ref{eq:quantizer} without
saturation at common scale $s$.  If quantization does not change any selected
path, the dequantized GL1F output
$\widehat y$ satisfies
\[
 |\widehat y-y|\le\frac{T+1}{2s}.
\]
For one output of a v2 model, replace $T$ by $R$.
\end{theorem}

\begin{proof}
Apply Lemma~\ref{lem:quant-error} to the base and every selected leaf, then use
the triangle inequality.
\end{proof}

\begin{lemma}[Sufficient split-margin condition]
\label{lem:split-margin}
Suppose neither $x$ nor threshold $\theta$ saturates.  If
\[
 |x-\theta|>\frac1s,
\]
then the truth value of $x>\theta$ equals that of
$Q_s(x)>Q_s(\theta)$.
\end{lemma}

\begin{proof}
If $x>\theta+1/s$, Lemma~\ref{lem:quant-error} implies
$Q_s(x)/s\ge x-1/(2s)>\theta+1/(2s)\ge Q_s(\theta)/s$.
The reverse ordering is analogous.
\end{proof}

\begin{corollary}[Argmax stability]
\label{cor:argmax}
Assume stable paths and unsaturated base/leaf quantization for a v2 model.  If
the difference between the largest and second-largest ideal logits exceeds
\[
 \frac{R+1}{s},
\]
the GL1F integer argmax returns the same class.  Equal integer logits are
resolved deterministically in favor of the lowest class index.
\end{corollary}

\begin{proof}
Theorem~\ref{thm:stable-path-error} perturbs each of the two logits by at most
$(R+1)/(2s)$.  Their order therefore cannot reverse under the stated margin.
The implementation replaces the incumbent only on strict improvement, which
establishes the tie rule.
\end{proof}

\begin{remark}[Why a global error bound needs a margin]
Near a split, quantization can select a different leaf.  The difference
between two leaves is not bounded by the rounding error alone.  Therefore
Theorem~\ref{thm:stable-path-error} must not be quoted without its stable-path
condition unless an additional bound on alternate-leaf values is supplied.
For production binary64 preprocessing, replace every $1/(2s)$ term above by
$(1/2+|\delta_s(z)|)/s$ for the corresponding quantized value, unless the
implementation is shown to emit the same integer as the abstract quantizer;
either alternative still requires that the corresponding int32 clamp not
saturate.  Provided neither operand saturates, a sufficient production split
margin is therefore
\[
 |x-\theta|>
 \frac{1+|\delta_s(x)|+|\delta_s(\theta)|}{s}.
\]
\end{remark}

\subsection{Limits of what the current registry proves}

For clarity, the live verifier names nine relations rather than using an
undefined ``applicable metadata'' category.  It checks
\[
\begin{split}
 &nFeatures=F,\qquad
 nTrees=\begin{cases}T&\text{(v1)},\\KR&\text{(v2)},\end{cases}
 \qquad depth=d,\\
 &baseQ=\begin{cases}a_0&\text{(v1)},\\0&\text{(v2 reserved-field convention)},\end{cases}
 \qquad scaleQ=s,\\
 &tablePtr=P,\qquad totalBytes=\ell,\qquad
 numChunks=N,\qquad chunkSize=c.
\end{split}
\]
The first five reconcile runtime registry metadata with the decoded core; the
last four reconcile the runtime and byte-information registry getters with
the storage manifest.

\begin{proposition}[No unconditional content binding]
\label{prop:no-binding}
For the deployed registry entry point, the existence of a record under
identifier $h$ does not imply
$h=\operatorname{keccak256}(B)$ for the bytes reconstructed from its pointer
table.
\end{proposition}

\begin{proof}
The registration predicate checks only that $h\ne0$ and that no prior record
exists under $h$; it accepts $h$ and the pointer metadata as independent
arguments.  Choose any nonzero $h$ and any accepted pointer tuple reconstructing
$B$ with $h\ne\operatorname{keccak256}(B)$.  The call can satisfy the remaining
fee, metadata, ToS, and licence checks.  Thus a valid transaction is a
counterexample to the implication.
\end{proof}

\begin{remark}[Executable assurance-boundary witness]
The isolated-chain integration test makes the counterexample concrete.  The
unmodified registry accepted this identifier:
\begin{center}
\small\ttfamily
0x60c713e19362bcb16b04fb6fc611821a\\
8a3fb90e1698490bf4819c2c44838875
\end{center}
for a pointer tuple whose 104 reconstructed bytes have hash
\begin{center}
\small\ttfamily
0xeb0ee33a46399a8592b6676035a7e580\\
c1f90e7761a2eb1f4a56d0f9f0e4ab18.
\end{center}
The runtime followed that record and returned integer prediction $1043$.
The production JavaScript chain loader now recomputes
$\operatorname{keccak256}(B)$ after exact reconstruction and rejects this
mismatch before returning model bytes.  This client-side check mitigates the
deployed assurance gap; it does not change the registry predicate.
\end{remark}

\begin{corollary}[Need for independent verification]
Theorem~\ref{thm:equivalence} applies to a registry entry only after an
independent verifier establishes Definitions~\ref{def:gl1f-v1} or
\ref{def:gl1f-v2}, Definition~\ref{def:manifest}, and all nine relations above.
If the registry key is additionally claimed as a content identifier, that
separate deployment-identity claim requires
$h=\operatorname{keccak256}(B)$; this equality is not a premise of execution
equivalence for the reconstructed bytes.
\end{corollary}

\begin{proposition}[Open model storage permits independent local inference]
\label{prop:public-model}
Under the documented GL1C storage design, an observer with chain-read access can
reconstruct the intentionally public model and evaluate it locally,
independently of the canonical runtime endpoint.
\end{proposition}

\begin{proof}
The pointer table and every chunk are public EVM runtime code retrievable by
address.  Definition~\ref{def:manifest} gives a constructive algorithm for
reassembling $B$.  Equations~\ref{eq:scalar-function} and
\ref{eq:vector-function} then define a local evaluator using only public
integers.  No secret required by evaluation is held by the fee-gated
contract.
\end{proof}

\begin{remark}
Mode-2 payment can still meter use of an official transaction endpoint,
produce an on-chain event, or support service-level economics.  It is not
model encryption, copy protection, or exclusive inference.
\end{remark}

\subsection{Verification requirements}

Applying the equivalence theorem to a concrete deployment requires an exact
core and storage reconstruction, canonical parser validation, agreement of
all nine registry/core/storage relations, and identical packed inputs.
A separately claimed content identifier requires recomputing the core's
Keccak-256 digest. Core bytes, chunk code, ordered addresses, registry
responses, packed vectors, and expected integer outputs allow these
conditions to be checked independently from a retained observation.

Source revision and compiler settings identify the build used in local
tests. Creation and deployed bytecode, together with their hashes, are
needed to examine any stronger source-to-deployment claim. Byte-identical
training results remain empirical evidence separate from the inference
theorem. These distinctions allow a failed reproduction to be assigned to a
specific representation, state, input, arithmetic, or implementation
premise.
